% --- Archon physics formalization source begin ---
% archon:physics
% archon:covers PhyXMiniProblems/problem_phyx_mini_0997.lean
% archon:source-report reports/phyx_mini/problem_phyx_mini_0997.source.json

\chapter{Physics problem phyx\_mini\_0997}
\label{ch:PhyXMiniProblems_problem_phyx_mini_0997}

\paragraph{Problem source.}
\#\# Physical scenario

An electric dipole is centered at the origin, with \textbackslash{}( \textbackslash{}vec\{p\} \textbackslash{}) in the direction of the \textbackslash{}( +y \textbackslash{})-axis.

\#\# Question

Derive an approximate expression for the electric field at a point \textbackslash{}( P \textbackslash{}) on the \textbackslash{}( y \textbackslash{})-axis for which \textbackslash{}( y \textbackslash{}) is much larger than \textbackslash{}( d \textbackslash{}). To do this, use the binomial expansion \textbackslash{}( (1 + x)\textasciicircum{}n \textbackslash{}approx 1 + nx + n(n - 1)x\textasciicircum{}2/2 + \textbackslash{}dots \textbackslash{}) (valid for the case \textbackslash{}( |x| < 1 \textbackslash{})).

\#\# Answer choices

- A: A: \textbackslash{}frac\{p\}\{4\textbackslash{}pi\textbackslash{}epsilon\_0 y\textasciicircum{}3\}

- B: B: \textbackslash{}frac\{p\}\{2\textbackslash{}pi\textbackslash{}epsilon\_0 y\textasciicircum{}3\}

- C: C: \textbackslash{}frac\{p\}\{2\textbackslash{}pi\textbackslash{}epsilon\_0 y\textasciicircum{}2\}

- D: D: \textbackslash{}frac\{2p\}\{2\textbackslash{}pi\textbackslash{}epsilon\_0 y\textasciicircum{}3\}

\#\# Figure caption (auxiliary; use the image as primary evidence)

The image depicts a vertical arrangement involving electric charges and vectors. Here’s a detailed description:

1. **Coordinate Axes**: 
   - The figure shows a vertical y-axis and a horizontal x-axis intersecting at a point labeled \textbackslash{}( O \textbackslash{}).

2. **Charges**:
   - A positive charge \textbackslash{}( +q \textbackslash{}) is located just above point \textbackslash{}( O \textbackslash{}).
   - A negative charge \textbackslash{}( -q \textbackslash{}) is positioned below point \textbackslash{}( O \textbackslash{}).
   - The distance between these two charges is labeled \textbackslash{}( d \textbackslash{}).

3. **Point \textbackslash{}( P \textbackslash{})**:
   - A point \textbackslash{}( P \textbackslash{}) is marked on the y-axis above both charges.
   - The distance from the center \textbackslash{}( O \textbackslash{}) to the point \textbackslash{}( P \textbackslash{}) is given as \textbackslash{}( y \textbackslash{}).

4. **Vectors**:
   - Vector \textbackslash{}( \textbackslash{}vec\{E\_+\} \textbackslash{}), pointing upwards from point \textbackslash{}( P \textbackslash{}), represents the electric field due to the positive charge.
   - Vector \textbackslash{}( \textbackslash{}vec\{E\_-\} \textbackslash{}), pointing downwards from point \textbackslash{}( P \textbackslash{}), represents the electric field due to the negative charge.
   - Vector \textbackslash{}( \textbackslash{}vec\{p\} \textbackslash{}), pointing upwards from \textbackslash{}( O \textbackslash{}), represents the dipole moment.

5. **Distances**:
   - The distance from \textbackslash{}( O \textbackslash{}) to the positive charge \textbackslash{}( +q \textbackslash{}) and from \textbackslash{}( O \textbackslash{}) to the negative charge \textbackslash{}( -q \textbackslash{}) is \textbackslash{}( d/2 \textbackslash{}).
   - The total distance from \textbackslash{}( P \textbackslash{}) to \textbackslash{}( O \textbackslash{}) is \textbackslash{}( y \textbackslash{}), therefore from \textbackslash{}( P \textbackslash{}) to the positive charge is \textbackslash{}( y - d/2 \textbackslash{}) and from \textbackslash{}( P \textbackslash{}) to the negative charge is \textbackslash{}( y + d/2 \textbackslash{}).
  
This diagram likely represents the electric field and dipole moment in a dipole system, with emphasis on the fields at point \textbackslash{}( P \textbackslash{}).

\#\# Dataset metadata

Electrostatics; Physical Model Grounding Reasoning, Numerical Reasoning

\paragraph{Recorded answer/context.}
B: B: \textbackslash{}frac\{p\}\{2\textbackslash{}pi\textbackslash{}epsilon\_0 y\textasciicircum{}3\}

\paragraph{Figure/image path.}
/root/proposal\_for\_physic/science-mango/phyx\_mini\_run/phyx\_data/test\_image/997.png

\paragraph{Formalization target.}
create a compiling Lean file with sorry bodies at `PhyXMiniProblems/problem\_phyx\_mini\_0997.lean`. The Lean declarations must preserve the physical quantities, dimensions or dimensional roles, figure labels, governing-law hypotheses, and final relation expressed by this problem.
Use Mathlib/Physlib names found through LeanExplore where available. If a physics API is missing, introduce faithful local abstractions rather than scalar placeholder aliases.

\begin{theorem}[Physics formalization target]
\label{thm:physics:phyx_mini_0997:target}
\lean{PhyXMiniProblems.ProblemPhyXMini0997.problem_phyx_mini_0997}
\uses{def:physics:phyx-mini-0997:phyxminiproblems-problemphyxmini0997-yaxisunit, def:physics:phyx-mini-0997:phyxminiproblems-problemphyxmini0997-observationpointp, def:physics:phyx-mini-0997:phyxminiproblems-problemphyxmini0997-axialelectricdipole, def:physics:phyx-mini-0997:phyxminiproblems-problemphyxmini0997-hasphysicalparameters, def:physics:phyx-mini-0997:phyxminiproblems-problemphyxmini0997-haschargeseparationdipolemoment, def:physics:phyx-mini-0997:phyxminiproblems-problemphyxmini0997-satisfiesaxialcoulombsuperposition, def:physics:phyx-mini-0997:phyxminiproblems-problemphyxmini0997-answeraxialfieldcomponent}
The assigned autoformalize agent should translate this physics problem into Lean declarations in the covered file, with theorem and lemma proof bodies written as `by sorry`.
\end{theorem}
\begin{proof}
This is an autoformalization task, not a proof task. Produce statements that can later be proved by the physics prover without weakening the physical model.
\end{proof}

% --- Archon declaration topology begin ---
\section{Declaration topology}
\begin{definition}[y Axis Unit]
  \label{def:physics:phyx-mini-0997:phyxminiproblems-problemphyxmini0997-yaxisunit}
  \lean{PhyXMiniProblems.ProblemPhyXMini0997.yAxisUnit}
  The unit vector in the positive y direction of the pictured coordinate system
\end{definition}
\begin{proof}
  This is the declared model definition of y Axis Unit.
\end{proof}
\begin{definition}[point On YAxis]
  \label{def:physics:phyx-mini-0997:phyxminiproblems-problemphyxmini0997-pointonyaxis}
  \lean{PhyXMiniProblems.ProblemPhyXMini0997.pointOnYAxis}
  A point with coordinate (0, y, 0) on the figure's y-axis
\end{definition}
\begin{proof}
  This is the declared model definition of point On YAxis.
\end{proof}
\begin{definition}[origin O]
  \label{def:physics:phyx-mini-0997:phyxminiproblems-problemphyxmini0997-origino}
  \lean{PhyXMiniProblems.ProblemPhyXMini0997.originO}
  \uses{def:physics:phyx-mini-0997:phyxminiproblems-problemphyxmini0997-pointonyaxis}
  The figure's origin O, which is also the center of the dipole
\end{definition}
\begin{proof}
  This is the declared model definition of origin O.
\end{proof}
\begin{definition}[observation Point P]
  \label{def:physics:phyx-mini-0997:phyxminiproblems-problemphyxmini0997-observationpointp}
  \lean{PhyXMiniProblems.ProblemPhyXMini0997.observationPointP}
  \uses{def:physics:phyx-mini-0997:phyxminiproblems-problemphyxmini0997-pointonyaxis}
  The figure's observation point P at axial coordinate y
\end{definition}
\begin{proof}
  This is the declared model definition of observation Point P.
\end{proof}
\begin{definition}[Axial Electric Dipole]
  \label{def:physics:phyx-mini-0997:phyxminiproblems-problemphyxmini0997-axialelectricdipole}
  \lean{PhyXMiniProblems.ProblemPhyXMini0997.AxialElectricDipole}
  Physical quantities and fields of the centered axial dipole. chargeMagnitude is the positive magnitude q, separation is the charge spacing d, and dipoleMoment is a vector readout in charge-times-length units. The three electric fields retain the distinct arrows E₊, E₋, and their total shown in the diagram
\end{definition}
\begin{proof}
  This is the declared model definition of Axial Electric Dipole.
\end{proof}
\begin{definition}[positive Charge Position]
  \label{def:physics:phyx-mini-0997:phyxminiproblems-problemphyxmini0997-positivechargeposition}
  \lean{PhyXMiniProblems.ProblemPhyXMini0997.positiveChargePosition}
  \uses{def:physics:phyx-mini-0997:phyxminiproblems-problemphyxmini0997-pointonyaxis, def:physics:phyx-mini-0997:phyxminiproblems-problemphyxmini0997-axialelectricdipole}
  The position (0, d / 2, 0) of the source charge +q
\end{definition}
\begin{proof}
  This is the declared model definition of positive Charge Position.
\end{proof}
\begin{definition}[negative Charge Position]
  \label{def:physics:phyx-mini-0997:phyxminiproblems-problemphyxmini0997-negativechargeposition}
  \lean{PhyXMiniProblems.ProblemPhyXMini0997.negativeChargePosition}
  \uses{def:physics:phyx-mini-0997:phyxminiproblems-problemphyxmini0997-pointonyaxis, def:physics:phyx-mini-0997:phyxminiproblems-problemphyxmini0997-axialelectricdipole}
  The position (0, -d / 2, 0) of the source charge -q
\end{definition}
\begin{proof}
  This is the declared model definition of negative Charge Position.
\end{proof}
\begin{definition}[distance PTo Positive]
  \label{def:physics:phyx-mini-0997:phyxminiproblems-problemphyxmini0997-distanceptopositive}
  \lean{PhyXMiniProblems.ProblemPhyXMini0997.distancePToPositive}
  \uses{def:physics:phyx-mini-0997:phyxminiproblems-problemphyxmini0997-axialelectricdipole}
  The figure-labelled distance y - d / 2 from P to +q
\end{definition}
\begin{proof}
  This is the declared model definition of distance PTo Positive.
\end{proof}
\begin{definition}[distance PTo Negative]
  \label{def:physics:phyx-mini-0997:phyxminiproblems-problemphyxmini0997-distanceptonegative}
  \lean{PhyXMiniProblems.ProblemPhyXMini0997.distancePToNegative}
  \uses{def:physics:phyx-mini-0997:phyxminiproblems-problemphyxmini0997-axialelectricdipole}
  The figure-labelled distance y + d / 2 from P to -q
\end{definition}
\begin{proof}
  This is the declared model definition of distance PTo Negative.
\end{proof}
\begin{definition}[Has Physical Parameters]
  \label{def:physics:phyx-mini-0997:phyxminiproblems-problemphyxmini0997-hasphysicalparameters}
  \lean{PhyXMiniProblems.ProblemPhyXMini0997.HasPhysicalParameters}
  \uses{def:physics:phyx-mini-0997:phyxminiproblems-problemphyxmini0997-axialelectricdipole}
  Positivity conditions for the vacuum and the scalar figure readouts q and d
\end{definition}
\begin{proof}
  This is the declared model definition of Has Physical Parameters.
\end{proof}
\begin{definition}[Has Charge Separation Dipole Moment]
  \label{def:physics:phyx-mini-0997:phyxminiproblems-problemphyxmini0997-haschargeseparationdipolemoment}
  \lean{PhyXMiniProblems.ProblemPhyXMini0997.HasChargeSeparationDipoleMoment}
  \uses{def:physics:phyx-mini-0997:phyxminiproblems-problemphyxmini0997-yaxisunit, def:physics:phyx-mini-0997:phyxminiproblems-problemphyxmini0997-axialelectricdipole}
  The defining dipole-moment law p⃗ = q d ŷ, including its +y direction
\end{definition}
\begin{proof}
  This is the declared model definition of Has Charge Separation Dipole Moment.
\end{proof}
\begin{definition}[Satisfies Axial Coulomb Superposition]
  \label{def:physics:phyx-mini-0997:phyxminiproblems-problemphyxmini0997-satisfiesaxialcoulombsuperposition}
  \lean{PhyXMiniProblems.ProblemPhyXMini0997.SatisfiesAxialCoulombSuperposition}
  \uses{def:physics:phyx-mini-0997:phyxminiproblems-problemphyxmini0997-yaxisunit, def:physics:phyx-mini-0997:phyxminiproblems-problemphyxmini0997-observationpointp, def:physics:phyx-mini-0997:phyxminiproblems-problemphyxmini0997-axialelectricdipole, def:physics:phyx-mini-0997:phyxminiproblems-problemphyxmini0997-distanceptopositive, def:physics:phyx-mini-0997:phyxminiproblems-problemphyxmini0997-distanceptonegative}
  Coulomb's law and linear superposition on the part of the y-axis above both charges. The positive source produces an upward field, the negative source a downward field, with the two figure distances y - d/2 and y + d/2
\end{definition}
\begin{proof}
  This is the declared model definition of Satisfies Axial Coulomb Superposition.
\end{proof}
\begin{definition}[binomial First Order]
  \label{def:physics:phyx-mini-0997:phyxminiproblems-problemphyxmini0997-binomialfirstorder}
  \lean{PhyXMiniProblems.ProblemPhyXMini0997.binomialFirstOrder}
  The first-order truncation 1 + n x of the binomial expansion
\end{definition}
\begin{proof}
  This is the declared model definition of binomial First Order.
\end{proof}
\begin{theorem}[inverse Square binomial first Order]
  \label{thm:physics:phyx-mini-0997:phyxminiproblems-problemphyxmini0997-inversesquare-binomial-firstorder}
  \lean{PhyXMiniProblems.ProblemPhyXMini0997.inverseSquare_binomial_firstOrder}
  \uses{def:physics:phyx-mini-0997:phyxminiproblems-problemphyxmini0997-binomialfirstorder}
  The requested binomial step for inverse-square factors: (1 + x)⁻² = 1 - 2x + o(x) as x → 0, within the stated domain |x| < 1
\end{theorem}
\begin{proof}
  The conclusion follows from the governing relations and figure readouts named in the statement of inverse Square binomial first Order.
\end{proof}
\begin{definition}[Answer Choice]
  \label{def:physics:phyx-mini-0997:phyxminiproblems-problemphyxmini0997-answerchoice}
  \lean{PhyXMiniProblems.ProblemPhyXMini0997.AnswerChoice}
  Labels of the four multiple-choice answers printed with the problem
\end{definition}
\begin{proof}
  This is the declared model definition of Answer Choice.
\end{proof}
\begin{definition}[answer Axial Field Component]
  \label{def:physics:phyx-mini-0997:phyxminiproblems-problemphyxmini0997-answeraxialfieldcomponent}
  \lean{PhyXMiniProblems.ProblemPhyXMini0997.answerAxialFieldComponent}
  \uses{def:physics:phyx-mini-0997:phyxminiproblems-problemphyxmini0997-answerchoice}
  Scalar +y electric-field component transcribed from each printed choice
\end{definition}
\begin{proof}
  This is the declared model definition of answer Axial Field Component.
\end{proof}
% --- Archon declaration topology end ---
% --- Archon physics formalization source end ---
